\documentclass[11pt,a4paper]{article}
\usepackage[margin=2.4cm]{geometry}
\usepackage{amsmath,amssymb,amsthm,mathtools}
\usepackage{booktabs}
\usepackage{enumitem}
\usepackage{xcolor}
\usepackage{hyperref}
\usepackage{fontspec}
\usepackage{polyglossia}
\setmainlanguage{russian}
\setotherlanguage{english}
\setmainfont{Times New Roman}
\setmonofont{Consolas}[Scale=0.90]

\hypersetup{
  colorlinks=true,
  linkcolor=blue!70!black,
  citecolor=blue!70!black,
  urlcolor=blue!70!black,
  pdftitle={Матричная кинематика ступенчатых передач на изотропном конусе: параболические жордановы замки SL2(Z) и бескорневое целочисленное сжатие Паде-Туэ в Lean 4},
  pdfauthor={Каратаев Дмитрий Александрович}
}

\theoremstyle{plain}
\newtheorem{theorem}{Теорема}[section]
\newtheorem{lemma}[theorem]{Лемма}
\newtheorem{proposition}[theorem]{Предложение}
\newtheorem{corollary}[theorem]{Следствие}
\theoremstyle{definition}
\newtheorem{definition}[theorem]{Определение}
\newtheorem{remark}[theorem]{Замечание}

\title{\textbf{Матричная кинематика ступенчатых передач на изотропном конусе:\\ параболические жордановы замки $\mathrm{SL}_2(\mathbb{Z})$ и бескорневое целочисленное сжатие Паде--Туэ--Вронского в Lean~4}}
\author{Д.\,А.~Каратаев\thanks{AI-Guiders Research (\url{https://github.com/AI-Guiders}). Формальная верификация в Lean~4 и аналитическая разработка выполнены при участии ИИ-партнёра \emph{Луч} (\emph{Luch}, AI-Guiders).}}
\date{}

\begin{document}
\maketitle

\begin{abstract}
Строится чисто алгебраический, бескорневой матричный метод и аппарат аппроксимаций Паде--Туэ для показательных диофантовых уравнений $A x^m - B y^n = C$ ($m,n\ge 2$) в кольце $\mathbb{Z}$ без использования иррациональных корней $\sqrt[3]{\widetilde{A}/\widetilde{B}}\in\mathbb{R}$ и аналитических оценок производных классического метода Туэ--Зигеля--Малера. Доказано, что каждая внутренняя пара показателей $(a,b)$ определяет $2\times 2$-матрицу $M_{a,b}$ определителя $C$, переход между которыми $\widetilde{P} = \frac{1}{C}M_{a,b}^{\mathrm{adj}}M_{a',b'}\in\mathrm{SL}_2(\mathbb{Q})$ допускает факторизацию ранга~$1$ и погружается в изотропный дискриминантный конус $X^2 + Y^2 - T^2 = \frac{1}{4}(\operatorname{Tr}\widetilde{P}^2 - 4)$. Вдоль линейных образующих конуса $X^2 - T^2 = 0$ ($\operatorname{Tr}\widetilde{P} = \mp 2$) зазор факторизуется в виде $C = (x^{m-a}\mp y^b)(x^a - y^b)$, а условие целочисленности $\widetilde{P}_k = I_2 \pm k N_0\in\mathrm{SL}_2(\mathbb{Z})$ ($N_0^2 = 0$) порождает квартет параболических жордановых замков $k\in\{6,12,15,40\}$ для $(3^2-2^3=1,\ 2^5-3^3=5,\ 5^3-3^4=44,\ 2^{10}-5^4=399)$.

Доказаны унимодулярная реконструкция оснований через матрицу показателей Фарея--Безу, трёхуровневый рекуррентный спуск $C_0 - (x^k+y^j)C_1 + x^k y^j C_2 = 0$ с квадратичным тождеством Ганкеля--Вронскиана $C_1^2 - C_0 C_2 = A B x^{m-2k}y^{n-2j}(x^k-y^j)^2$, невырожденность башни вронскианов Паде--Туэ $\det\mathcal{M}_N = -(3N+1)X Y C^{2N+1}$, межэтажный вронскиан $\det\mathcal{W}_{2,3} = (U-V)^3(U^2+3UV+V^2)$ (\texttt{thm11}) и бескорневое тождество кубического сжатия Паде $U I_1 \Phi_3(a\widehat{P}_1, b\widehat{Q}_1) = a^3(U-V)^3(U+V) - (U b^3 - V a^3)\widehat{Q}_1^3$ в $d^2=9$ клетках вычетов (\texttt{thm10}). Все теоремы (\texttt{thm6}, \texttt{thm7}, \texttt{thm10}, \texttt{thm11}) формально верифицированы в ядре \textbf{Lean~4} (\texttt{0 sorryAx}, \texttt{0 Classical.choice}).

\medskip\noindent
\textbf{Ключевые слова:} показательные диофантовы уравнения, уравнение Пиллаи, аппроксимации Паде--Туэ, определитель Вронского, изотропный дискриминантный конус, параболические жордановы матрицы $\mathrm{SL}_2(\mathbb{Z})$, Lean~4.
\end{abstract}

\section{Матричная декомпозиция на изотропном дискриминантном конусе и квартет параболических жордановых замков \texorpdfstring{$\mathrm{SL}_2(\mathbb{Z})$}{SL2(Z)}}
\label{sec:isotropic-cone}

\begin{theorem}[Матричная декомпозиция, образующие изотропного дискриминантного конуса, факторизация $C = (x^{m-a}\mp y^b)(x^a-y^b)$ и параболический жорданов квартет $\mathrm{SL}_2(\mathbb{Z})$ $k\in\{6,12,15,40\}$]
\label{thm:main-6}
Для показательного диофантова уравнения $A x^m - B y^n = C$ ($m,n\ge 2$) выполнены следующие матричные соотношения (\texttt{thm6\_pillai\_skeletal\_matrix\_decomposition\_and\_parabolic\_sl2\_jordan\_triptych}):
\begin{enumerate}[label=\textbf{(\roman*)},leftmargin=*]
\item каждая внутренняя пара показателей $(a,b)$ определяет матрицу $M_{a,b} = \begin{pmatrix} A x^{m-a} & B y^{n-b} \\ y^b & x^a \end{pmatrix}$ с определителем $\det M_{a,b} = C$, а нормированный переход $\widetilde{P} = \frac{1}{C} M_{a,b}^{\mathrm{adj}} M_{a',b'}\in\mathrm{SL}_2(\mathbb{Q})$ ($\det\widetilde{P} = +1$, $\Delta a = a'-a$, $\Delta b = b'-b$) удовлетворяет факторизации ранга~$1$:
  \[
    \widetilde{P} \;-\; \begin{pmatrix} y^{\Delta b} & 0 \\ 0 & x^{\Delta a} \end{pmatrix}
    \;=\; \frac{x^{\Delta a}y^{\Delta b}-1}{C}\,\begin{pmatrix} -x^a \\ y^b \end{pmatrix}\begin{pmatrix} A x^{m-a'} & B y^{n-b'} \end{pmatrix};
  \]
\item в координатах квадрики $X = (\widetilde{P}_{11}-\widetilde{P}_{22})/2$, $Y = (\widetilde{P}_{12}+\widetilde{P}_{21})/2$, $T = (\widetilde{P}_{12}-\widetilde{P}_{21})/2$ дискриминант равен $\frac{1}{4}(\operatorname{Tr}\widetilde{P}^2 - 4) = X^2 + Y^2 - T^2$. На симметричных разбиениях ($m=a+a'$, $n=b+b'$, $A=B=1$) $Y\equiv 0$, образующая $X - T$ удовлетворяет $2C(X - T) = (x^a - y^b)^2 - (x^{a'} - y^{b'})^2$, а на изотропном конусе $X^2 - T^2 = 0$ ($\operatorname{Tr}\widetilde{P} = \mp 2$) зазор факторизуется в виде $C = (x^{m-a}\mp y^b)(x^a - y^b)$, порождая при $\widetilde{P}_k\in\mathrm{SL}_2(\mathbb{Z})$ четыре канонических жордановых замка $k\in\{6,12,15,40\}$.
\end{enumerate}
\end{theorem}

\section{Трёхуровневый спуск Ганкеля--Вронскиана, башня Паде--Туэ и бескорневое кубическое сжатие}
\label{sec:pade-thue}

\begin{theorem}[Унимодулярная матрица Фарея--Безу, трёхуровневый рекуррентный спуск, моноид башни Паде--Туэ, невырожденность соседних этажей $|I_{\mathcal{R}}| + |I_{\mathcal{R}\mathbin{+\!+}[r]}| \ge 1$ и двусторонняя вилка Паде, \texttt{thm7}]
\label{thm:main-7}
Для взаимно простых показателей $\gcd(m,n)=1$ уравнения $A x^m - B y^n = C$ родительские узлы медианты Фарея--Безу $(a,b)$ и $(a',b') = (m-a,n-b)$ образуют $E_{\mathrm{B\acute{e}zout}} = \begin{pmatrix} a & b \\ a' & b' \end{pmatrix}\in\mathrm{SL}_2(\mathbb{Z})$ ($a b' - a' b = +1$), восстанавливая основания: $(x^a)^{b'} = x\cdot (x^{a'})^b$ и $(y^{b'})^a = y\cdot (y^b)^{a'}$. Вдоль любого луча $(k,j)$ последовательные разности $C_t = A x^{m-tk} - B y^{n-tj}$ удовлетворяют $C_0 - (x^k+y^j)C_1 + x^k y^j C_2 = 0$ и квадратичному тождеству Ганкеля--Вронскиана $C_1^2 - C_0 C_2 = A B x^{m-2k}y^{n-2j}(x^k-y^j)^2$. На счётном моноиде $\mathbb{Z}^{<\infty}$ $M$-этажных передаточных каскадов Паде--Туэ $\mathcal{R} = [r_0,r_1,\dots,r_{M-1}]$ (\texttt{padeTowerP}, \texttt{padeTowerQ}, \texttt{padeTowerProdPlus}, \texttt{padeTowerProdMinus}) композиция удовлетворяет законам $P_{\mathcal{R}_1\mathbin{+\!+}\mathcal{R}_2}(C,B) = P_{\mathcal{R}_2}(P_{\mathcal{R}_1},Q_{\mathcal{R}_1})$, $Q_{\mathcal{R}_1\mathbin{+\!+}\mathcal{R}_2}(C,B) = Q_{\mathcal{R}_2}(P_{\mathcal{R}_1},Q_{\mathcal{R}_1})$, $\Pi_\pm(\mathcal{R}_1\mathbin{+\!+}\mathcal{R}_2) = \Pi_\pm(\mathcal{R}_1)\Pi_\pm(\mathcal{R}_2)$ ($\Pi_+(\mathcal{R})=\prod_{m=0}^{M-1}r_m$, $\Pi_-(\mathcal{R})=\prod_{m=0}^{M-1}(2-r_m)$), $P_{\mathcal{R}}+Q_{\mathcal{R}} = \Pi_+(\mathcal{R})(C+B)$, $P_{\mathcal{R}}-Q_{\mathcal{R}} = \Pi_-(\mathcal{R})A$ ($A=C-B$) и $Q_{\mathcal{R}}(u-v) - I_{\mathcal{R}} = -u\,\Pi_-(\mathcal{R})A$ для $I_{\mathcal{R}} := u P_{\mathcal{R}} - v Q_{\mathcal{R}}$. Более того, между любым этажом $\mathcal{R}$ и следующим этажом $\mathcal{R}' = \mathcal{R}\mathbin{+\!+}[k+1]$ ($P_{\mathcal{R}'} = P_{\mathcal{R}} + k Q_{\mathcal{R}}$, $Q_{\mathcal{R}'} = k P_{\mathcal{R}} + Q_{\mathcal{R}}$) выполнены межэтажные тождества Вронского:
\begin{equation}
  \label{eq:adjacent-pade-wronskian}
  \begin{aligned}
    I_{\mathcal{R}'} - (k+1)I_{\mathcal{R}} &= -k(u+v)\,\Pi_-(\mathcal{R})\,A, &\quad
    I_{\mathcal{R}'} + (k-1)I_{\mathcal{R}} &= k(u-v)\,\Pi_+(\mathcal{R})(C+B), \\
    Q_{\mathcal{R}'}I_{\mathcal{R}} - Q_{\mathcal{R}}I_{\mathcal{R}'} &= k\,u\,\Pi_-(\mathcal{R})\,\Pi_+(\mathcal{R})\,A(C+B), &\quad
    P_{\mathcal{R}'}Q_{\mathcal{R}} - Q_{\mathcal{R}'}P_{\mathcal{R}} &= -k\,\Pi_-(\mathcal{R})\,\Pi_+(\mathcal{R})\,A(C+B).
  \end{aligned}
\end{equation}
Следовательно, при $k(u+v)\,\Pi_-(\mathcal{R})A \ne 0$ две соседние невязки столкновения никогда не обращаются в нуль одновременно:
\[
  \neg\bigl(I_{\mathcal{R}} = 0 \;\land\; I_{\mathcal{R}'} = 0\bigr), \qquad
  |I_{\mathcal{R}}| + |I_{\mathcal{R}'}| \;\ge\; 1 \quad\bigl(\max(|I_{\mathcal{R}}|,|I_{\mathcal{R}'}|) \ge 1\bigr),
\]
а противоположные знаки соседних этажей ($I_{\mathcal{R}} \le 0 \le I_{\mathcal{R}'}$ или $I_{\mathcal{R}'} \le 0 \le I_{\mathcal{R}}$) зажимают пару $(u,v)$ в двустороннюю целочисленную вилку $u P_{\mathcal{R}} \le v Q_{\mathcal{R}} \land v Q_{\mathcal{R}'} \le u P_{\mathcal{R}'}$ (соответственно $u P_{\mathcal{R}'} \le v Q_{\mathcal{R}'} \land v Q_{\mathcal{R}} \le u P_{\mathcal{R}}$).
\end{theorem}

\begin{theorem}[Бескорневое тождество кубического сжатия Паде в $9$ клетках вычетов, \texttt{thm10}]
\label{thm:main-10}
Для любых $(U,V,a,b)\in\mathbb{Z}^4$ инвариант $I_1 = a\widehat{P}_1 - b\widehat{Q}_1$ ($\widehat{P}_1 = U+2V,\ \widehat{Q}_1 = 2U+V$) и кубический фактор $\Phi_3(x,y) = x^2 + xy + y^2$ удовлетворяют тождеству
\begin{equation}
  \label{eq:root-free-cubic-compression}
  U\,I_1\,\Phi_3\bigl(a\widehat{P}_1,\,b\widehat{Q}_1\bigr) \;\equiv\; a^3(U-V)^3(U+V) \;-\; \bigl(U b^3 - V a^3\bigr)\,\widehat{Q}_1^3.
\end{equation}
\end{theorem}

\begin{theorem}[Тождество межэтажного Вронскиана между этажами $k=2$ и $k=3$, \texttt{thm11}]
\label{thm:main-11}
Для любых $(U,V)\in\mathbb{Z}^2$ межэтажный вронскиан между квадратичной парой Паде ($\widehat{P}_{1,2}=U+3V,\ \widehat{Q}_{1,2}=3U+V$) и кубической парой Паде ($\widehat{P}_{1,3}=U+2V,\ \widehat{Q}_{1,3}=2U+V$) факторизуется тождественно:
\begin{equation}
  \label{eq:cross-floor-wronskian}
  \det\mathcal{W}_{2,3} \;=\; \widehat{Q}_{1,2}^2\widehat{P}_{1,3}^3 - \widehat{P}_{1,2}^2\widehat{Q}_{1,3}^3 \;\equiv\; (U - V)^3\bigl(U^2 + 3UV + V^2\bigr).
\end{equation}
\end{theorem}

\section{Формальная верификация в Lean~4}
Формальное доказательство Теорем~\ref{thm:main-6} и~\ref{thm:main-7} (\texttt{thm6}, \texttt{thm7}, \texttt{thm10}, \texttt{thm11}) содержится в модуле \texttt{PillaiDiophantine.lean} в репозитории \url{https://github.com/AI-Guiders/scientific-papers} (\texttt{0 sorryAx}, \texttt{0 Classical.choice}).

\begin{thebibliography}{9}
\bibitem{Pillai1931}
S.\,S.~Pillai, ``On the inequality $0 < a^x - b^y \le n$,'' \emph{J.~Indian Math.~Soc.} \textbf{19} (1931), 1--11.
\bibitem{Thue1909}
A.~Thue, ``\"Uber Ann\"aherungswerte algebraischer Zahlen,'' \emph{J.~reine angew.~Math.} \textbf{135} (1909), 284--305.
\bibitem{Moura2021}
L.~de~Moura and S.~Ullrich, ``The Lean~4 theorem prover and programming language,'' in \emph{CADE~28}, LNCS vol.~12699, Springer, 2021, pp.~625--635.
\end{thebibliography}

\end{document}
